% ============================================================
%  Section 6: Experimental Setup
%  H&E to Sirius Red Virtual Staining Paper
% ============================================================

\section{Experimental Setup}
\label{sec:exp_setup}

% -------------------------------------------------------
\subsection{Image Preprocessing and Tiling}
\label{sec:tiling}
% -------------------------------------------------------

Whole-slide images~(WSIs) were tiled into non-overlapping $256{\times}256$\,pixel
patches using a stride equal to the tile size.
Tiles with less than 50\% tissue content (assessed against the corresponding tissue
mask) were discarded from the training set; no tissue filtering was applied to the test
set.
All tiles were resized to $256{\times}256$ pixels at extraction time and normalised
channel-wise to $[-1,\,1]$ by subtracting $0.5$ and dividing by $0.5$ per channel.
The resulting tiles are organised into numbered per-WSI subdirectories
(\texttt{001/}, \texttt{002/}, \ldots) for domain~A~(H\&E) and domain~B~(SR) separately,
enabling reproducible data-fraction subsets via contiguous folder ranges
(Section~\ref{sec:scaling}).

% -------------------------------------------------------
\subsection{Training Protocol}
\label{sec:training_protocol}
% -------------------------------------------------------

All models are trained with the Adam optimiser~\cite{kingma2015adam} using a learning
rate of $2{\times}10^{-4}$, $\beta_1{=}0.5$, and $\beta_2{=}0.999$, with a batch size
of~1.
Training is \textit{step-based} rather than epoch-based: each optimiser update
constitutes one step, ensuring a fair comparison across data-fraction configurations
that differ in the number of available tiles.
All configurations use automatic mixed precision~(AMP) to reduce GPU memory and wall
time without loss of numerical stability.

Single-stage models~(CycleGAN, UNIT, MUNIT, DCLGAN) are trained for
$5{\times}10^6$~steps.
Multi-stage models use the same total budget distributed across their stages:

\begin{itemize}
  \item \textbf{UVCGAN:} $1.25{\times}10^6$~steps of masked image pretraining
    (Stage~1) followed by $3.75{\times}10^6$~steps of cycle-consistent finetuning
    (Stage~2), with Stage~2 initialised from the Stage~1 checkpoint.

  \item \textbf{MIUDiff:} $3.5{\times}10^6$~steps of unconditional DDPM pretraining
    on domain~B (Stage~1), $7.5{\times}10^5$~steps of conditional MI-guided finetuning
    (Stage~2), and $7.5{\times}10^5$~steps of patch-contrastive finetuning (Stage~3),
    with each stage initialised from the previous stage's final checkpoint.
\end{itemize}

Model checkpoints are saved every $2.5{\times}10^5$~steps; generator losses are logged
every $10^3$~steps to \texttt{loss\_log.csv} for post-hoc convergence analysis.
All experiments are run on a single GPU~(\texttt{paula} partition) with 8~CPU cores and
32\,GB RAM per job.
Training is fully resumable: if a run is interrupted, it automatically restores the
latest available checkpoint and continues from that step.

% -------------------------------------------------------
\subsection{Scaling Study Design}
\label{sec:scaling}
% -------------------------------------------------------

We conduct a full factorial experiment over three axes:

\paragraph{Model architecture.}
The six models evaluated are CycleGAN, UNIT, MUNIT, DCLGAN, UVCGAN, and MIUDiff,
covering cycle-consistent GANs, shared-latent-space GANs, content--style
disentanglement, contrastive learning, hybrid ViT--GANs, and score-based diffusion
(detailed in Section~\ref{sec:models}).

\paragraph{Generator capacity.}
Each model is trained at three generator sizes --- Small~(${\approx}10$\,M), Medium~(${\approx}50$\,M), and Large~(${\approx}100$\,M) parameters --- by adjusting the base channel width
$n_{gf}$ and depth $n_{\text{blocks}}$ (or equivalent per model family; see
Table~\ref{tab:model_sizes}).
All size variants within a model family share the same training protocol, data, and
random seed, so that any performance difference is attributable solely to generator
capacity.

\paragraph{Training data fraction.}
The dataset is partitioned at the WSI level into three nested subsets: 25\%, 50\%, and
100\% of the full training set, corresponding to WSI folder ranges
$\llbracket 1,7 \rrbracket$, $\llbracket 1,15 \rrbracket$, and $\llbracket 1,30 \rrbracket$
respectively.
Smaller subsets are strict subsets of larger ones, so differences across data fractions
are not confounded by variation in which slides are included.
The same fixed test set~(held-out WSIs not included in any training fraction) is used
for evaluation across all configurations.

\vspace{0.5em}
\noindent
The full factorial yields $6\text{ models} \times 3\text{ sizes} \times 3\text{ data
fractions} = \mathbf{54}$ experimental configurations, each trained and evaluated
independently as a separate SLURM array job.

% -------------------------------------------------------
\subsection{Sirius Red Collagen Segmentation Model}
\label{sec:sr_seg_model}
% -------------------------------------------------------

The task-based evaluations~(Sections~\ref{sec:task_eval} and~\ref{sec:cross_stain})
require a frozen collagen segmentation model applied identically to real and virtual
SR images.
We adopt an off-the-shelf nnU-Net~v2 model~\cite{[TODO: nnUNet SR publication]} trained
for PSR-positive area segmentation on Sirius Red WSIs; performance figures for this model
are reported in the original publication.
Because the same model is applied without modification to all conditions, any systematic
segmentation errors cancel out in the comparison: residual differences in predicted
collagen area reflect differences in staining fidelity between generative models rather
than segmentation imprecision.

% -------------------------------------------------------
\subsection{Inference Configuration}
\label{sec:inference}
% -------------------------------------------------------

At test time, all GAN-based models perform a single forward pass through the trained
$G_{A \to B}$ generator per tile.
MIUDiff uses DDIM sampling~\cite{song2021ddim} over $T_{\text{sample}}{=}300$
deterministic reverse-diffusion steps with classifier-free MI guidance at scale
$\eta{=}1.0$.
For all models, generated tiles are denormalised from $[-1,\,1]$ to $[0,\,1]$ and
saved as TIFF files, preserving the original tile identifiers for downstream metric
computation.

% -------------------------------------------------------
\subsection{Reproducibility}
\label{sec:reproducibility}
% -------------------------------------------------------

All 54 training jobs are defined in a single SLURM job-array script that encodes every
architectural hyperparameter and data-fraction argument explicitly, ensuring full
reproducibility from a single command.
Each checkpoint stores the complete model configuration alongside the optimiser state
and accumulated training time, so any run can be resumed or re-evaluated at any saved
step.
The dataset, pre-trained model checkpoints, and training scripts will be released
publicly alongside this paper.


% ============================================================
%  New bibliography entries for this section
%  (merge with references.bib)
% ============================================================
%
% @inproceedings{kingma2015adam,
%   author    = {Kingma, Diederik P. and Ba, Jimmy},
%   title     = {Adam: A Method for Stochastic Optimization},
%   booktitle = {ICLR},
%   year      = {2015},
%   eprint    = {1412.6980}
% }
%
% (Shared entries already listed:
%  song2021ddim)
